% ===== sec/04_trace_language.tex =====
\section{The Trace Language of an Agent}
\label{sec:trace}

This section is the formal core of the paper.  We define an agent
as a tuple of seven objects, three of which are derivative.  We
then define the agent's trace language and prove that the
definition satisfies the Chomsky hierarchy verbatim.  The
construction is austere and self-contained; nothing beyond
Section~\ref{sec:preliminaries} is assumed.

\subsection{Operation alphabet}

\begin{definition}[Operation alphabet]
\label{def:opalpha}
An \emph{operation alphabet} is a finite, non-empty set $\Sigma$
whose elements $\sigma\in\Sigma$ each denote an effectful
computation \emph{name}, called an \emph{operation symbol}.
Each $\sigma$ is associated with a (possibly partial) computable
function $\llbracket\sigma\rrbracket : X \rightharpoonup Y$
acting on a state space $X$ and producing a successor state in
$Y\subseteq X$, so that the abstract symbol carries an effect when
executed.  The function $\llbracket\cdot\rrbracket$ is fixed once
and for all by the runtime environment and is not a property of
the agent.
\end{definition}

\noindent For the empirical instantiation in
Section~\ref{sec:expsetup}, $\Sigma$ is the union of the column
entries of \texttt{trace\_language\_kaggle.csv}.  Examples include
\texttt{load\_vector\_db(gemini\_embeddings.json)},
\texttt{compose\_planner\_prompt},
\texttt{invoke\_gemini(gemini-2.5-flash-lite)},
\texttt{run\_phase\_unit\_tests(phase; comp\_dir)}, and
\texttt{train\_stacked\_ensemble(LGBM+XGB+HistGB+linear\_meta)}.
The operation symbol is the \emph{name plus its parenthesised
argument list as written in the source}; two invocations of the
same name with different arguments are different operation
symbols, exactly as Liskov's abstract data types prescribe for
parameterised operations~\cite{liskov1974abstraction}.

\subsection{Agent}

\begin{definition}[Agent]
\label{def:agent}
An \emph{agent} is a 7-tuple
\[
\mathcal{A} \;=\; \bigl(\Sigma_\mathcal{A},\;
                       Q_\mathcal{A},\;
                       q_0^\mathcal{A},\;
                       X_\mathcal{A},\;
                       \delta_\mathcal{A},\;
                       F_\mathcal{A},\;
                       h_\mathcal{A}\bigr)
\]
where
\begin{itemize}
\item $\Sigma_\mathcal{A} \subseteq \Sigma$ is the agent's
operation alphabet;
\item $Q_\mathcal{A}$ is a (possibly countably infinite)
\emph{configuration space};
\item $q_0^\mathcal{A} \in Q_\mathcal{A}$ is the \emph{initial
configuration};
\item $X_\mathcal{A}$ is the (computable) \emph{input domain};
\item $\delta_\mathcal{A} : Q_\mathcal{A} \times X_\mathcal{A}
\rightharpoonup Q_\mathcal{A} \times \Sigma_\mathcal{A}^*$ is the
agent's \emph{step relation}: at configuration $q$ on input $x$
it returns a successor configuration $q'$ together with a finite
string of operation symbols emitted by the step;
\item $F_\mathcal{A} \subseteq Q_\mathcal{A}$ is the \emph{halting
set};
\item $h_\mathcal{A}: F_\mathcal{A} \to \{\textsf{accept},
\textsf{reject}\}$ is the \emph{halting verdict}.
\end{itemize}

The objects $X_\mathcal{A}, \delta_\mathcal{A}, h_\mathcal{A}$
are derivative: $X_\mathcal{A}$ is determined by $\Sigma$ and the
runtime environment; $\delta_\mathcal{A}$ is the agent's
\emph{program}; and $h_\mathcal{A}$ is fixed by the convention
that a configuration in $F_\mathcal{A}$ accepts iff it carries
the post-condition for which the trace was emitted.
\end{definition}

\begin{remark}
The 7-tuple is the minimal one that lets us speak of an emitted
\emph{string} per step.  A 5-tuple of the
$(Q,\Sigma,\delta,q_0,F)$ Hopcroft--Ullman~\cite{hopcroft2006automata}
form would force $\delta$ to emit a single symbol per step, which
is too restrictive: a single ``call to the planner'' in our
empirical agent emits a structured prompt-build sequence that we
prefer to record as a string of three or four operation symbols.
The choice of $\Sigma_\mathcal{A}^*$ rather than $\Sigma_\mathcal{A}$
in the codomain of $\delta_\mathcal{A}$ is therefore essential.
\end{remark}

\subsection{Run, trace, and trace language}

\begin{definition}[Run]
\label{def:run}
A \emph{run} of $\mathcal{A}$ on input $x\in X_\mathcal{A}$ is a
finite or infinite sequence
\[
\rho = \bigl((q_0,\, w_0),\; (q_1,\, w_1),\; (q_2,\, w_2),\;
\dots\bigr)
\]
with $q_0 = q_0^\mathcal{A}$, $w_0 = \varepsilon$, and for each
$i\ge 1$, $\delta_\mathcal{A}(q_{i-1}, x) = (q_i, w_i)$ defined.
The run is \emph{halting} if it is finite and its last
configuration $q_n$ lies in $F_\mathcal{A}$; it is
\emph{accepting} if additionally $h_\mathcal{A}(q_n) =
\textsf{accept}$.
\end{definition}

\begin{definition}[Trace]
\label{def:trace}
The \emph{trace} of a run $\rho = ((q_0,w_0),
\ldots,(q_n,w_n))$ is the concatenation
\[
\tau(\rho) \;=\; w_0 \cdot w_1 \cdots w_n \;\in\;
\Sigma_\mathcal{A}^*.
\]
We write $\tau(\mathcal{A}, x)$ for the trace of the run of
$\mathcal{A}$ on $x$ when that run is unique.
\end{definition}

\begin{definition}[Trace language]
\label{def:tracelang}
The \emph{trace language} of $\mathcal{A}$ is
\[
L(\mathcal{A}) \;=\;
\bigl\{\,\tau(\rho) \;:\; \rho \text{ is an accepting run of }
\mathcal{A} \text{ on some } x\in X_\mathcal{A}\,\bigr\}
\;\subseteq\; \Sigma_\mathcal{A}^*.
\]
\end{definition}

\subsection{Trace-equivalence}

\begin{definition}[Trace equivalence]
\label{def:traceeq}
Two agents $\mathcal{A}_1, \mathcal{A}_2$ over a common alphabet
$\Sigma_\mathcal{A}$ are \emph{trace-equivalent}, written
$\mathcal{A}_1 \equiv_\tau \mathcal{A}_2$, iff
$L(\mathcal{A}_1) = L(\mathcal{A}_2)$.
\end{definition}

\noindent The austerity claim of the paper, in two lines, is:
\emph{An agent is its trace language, and only its trace
language.}  Trace-equivalent agents are the same agent.  The
internal architecture (LLM, neural network, hand-coded controller,
mixture of all three) is irrelevant to the identity criterion.

\subsection{Multi-agent systems}

\begin{definition}[Multi-agent system]
\label{def:multi}
A \emph{multi-agent system} is a finite tuple $\mathcal{M} =
(\mathcal{A}_1, \dots, \mathcal{A}_k, R)$ where each
$\mathcal{A}_i$ is an agent and $R \subseteq \Sigma^*\times\{1,
\dots,k\}\times\{1,\dots,k\}$ is an \emph{orchestration relation}
specifying which agent's emission, on which prefix, may be
consumed by which other agent.
\end{definition}

The trace language of $\mathcal{M}$ is the set of strings over the
common alphabet $\Sigma = \bigcup_i \Sigma_{\mathcal{A}_i}$ that
are obtainable by the orchestrated shuffle: pick an accepting run
of each $\mathcal{A}_i$, then interleave their traces in any order
consistent with $R$.  Formally,
\[
L(\mathcal{M}) \;=\;
\bigcup_{(\tau_1, \dots, \tau_k)\,\in\,
        L(\mathcal{A}_1)\times\cdots\times L(\mathcal{A}_k)}
\textsf{shuffle}_R(\tau_1, \dots, \tau_k)
\]
where $\textsf{shuffle}_R$ is the standard interleaving
restricted by $R$~\cite{hopcroft2006automata}.

\subsection{The CSV files \emph{are} languages}

We make explicit the empirical correspondence used throughout the
sequel.  The header row of \texttt{trace\_language\_kaggle.csv}
names the sub-agents:
\[
\mathcal{A}_1 = \texttt{research\_agent},\quad
\dots,\quad
\mathcal{A}_{13} = \texttt{kaggle\_trainer}.
\]
For $1\le i\le 13$, the $i$-th column of the file (skipping empty
cells) is a single string $\tau_i \in \Sigma_{\mathcal{A}_i}^*$.
This string is one accepting run of $\mathcal{A}_i$ on the
\textsc{Spaceship Titanic} input.  The trace language
$L(\mathcal{A}_i)$ contains, in principle, infinitely many other
strings; the column records exactly the one that the
recursive enumeration emitted.

The same correspondence holds for
\texttt{trace\_language\_original.csv}, with the difference that
its header has only twelve columns (the start language has no
\texttt{kaggle\_trainer} agent) and its row entries are
substantially shorter.

\subsection{Closure under intersection with the verifier}

We now state the closure result that is the architectural pivot
of the paper.

\begin{theorem}[Verifier closure]
\label{thm:closure}
Let $\mathcal{G}$ be an agent with $L(\mathcal{G})$ in Chomsky
class $\mathcal{L}_i$ for some $i\in\{0,1,2,3\}$.  Let
$\mathcal{V}$ be an agent with $L(\mathcal{V}) \in \mathcal{L}_3$
(a regular language).  Define the \emph{verified composition}
$\mathcal{G}\,\Vert\,\mathcal{V}$ as the agent that emits a
trace $\tau$ iff $\tau \in L(\mathcal{G})$ and $\tau$ is accepted
by $\mathcal{V}$ when treated as input to its DFA.  Then
\[
L(\mathcal{G}\,\Vert\,\mathcal{V})
\;=\; L(\mathcal{G}) \cap L(\mathcal{V})
\;\in\; \mathcal{L}_i.
\]
In particular, the Chomsky class of the verified trace language
is no greater than the Chomsky class of $\mathcal{G}$ alone, and
the membership problem for $L(\mathcal{G}\,\Vert\,\mathcal{V})$
is decidable in linear time whenever a candidate trace is
presented to $\mathcal{V}$'s DFA.
\end{theorem}

\begin{proof}
Each Chomsky class $\mathcal{L}_i$ is closed under intersection
with regular languages.  For $i = 3$, this is the standard
product-construction~\cite{rabin1959finite}.  For $i = 2$, it is
the standard product of a PDA with a DFA, which is again a
PDA~\cite{hopcroft2006automata}.  For $i = 1$, it follows from
LBA closure under intersection (regular languages are
context-sensitive).  For $i = 0$, it is the standard product
of a Turing machine with a DFA, which is again a Turing machine
\cite{turing1936computable,hopcroft2006automata}.

The decidability claim follows because $\mathcal{V}$ is a DFA:
given a candidate trace $\tau$ of length $n$, the DFA decides
$\tau \in L(\mathcal{V})$ in $O(n)$ time and constant space.
\qed
\end{proof}

\begin{corollary}[Verification budget]
\label{cor:budget}
Verification of the verified trace language
$L(\mathcal{G}\,\Vert\,\mathcal{V})$ on a candidate trace $\tau$
of length $n$ runs in $O(n)$ time and $O(1)$ space.  The Type-0
class of $\mathcal{G}$ does not enlarge the verifier's complexity.
\end{corollary}

\subsection{The empirical signature of a Type-0 trace language}

A trace language $L(\mathcal{A})$ is Type-0 iff $\mathcal{A}$ can
emit traces whose length grows without a recursive bound in the
input size, and whose contents reflect arbitrary recursive
enumerations of the operation alphabet.  In our empirical
setting, $\mathcal{A}$'s Type-0 status is witnessed by the
\texttt{kaggle\_trainer} column of
\texttt{trace\_language\_kaggle.csv}, which contains a strictly
longer sequence (36 non-empty entries) than any column of
\texttt{trace\_language\_original.csv} (which contains at most
fourteen).  The lengthening is the operational footprint of the
recursive enumeration: the start language emitted twelve seed
agents; the binding to the Type-3 verifier forced enumeration of
arbitrarily long check-train-validate cycles; the result was a
strictly more complex trace, materialised on disk by the
recursive Type-0 step.

The undecidability that Robinson's
theorem~\cite{robinson1949definability,davis1961decision} fixes
on the Type-0 class is therefore a real property of $\mathcal{A}$:
were it not for the Type-3 verifier, no algorithm would decide
whether a candidate trace lay in $L(\mathcal{A})$.  The verifier
is the entire reason the system is decidable in practice.
